\documentclass[10pt,a4paper]{amsart}
\usepackage[margin=19mm]{geometry}
\usepackage{amsmath,amssymb,mathtools}
\usepackage[scr=rsfs,cal=euler]{mathalfa}
\usepackage{xfrac}
\usepackage[unicode,hidelinks,bookmarks=false]{hyperref}
\newcommand{\ie}{i.e.\ }
\newcommand{\cf}{cf.\ }
\newcommand{\resp}{respectively\ }
\newcommand{\Subsection}{\subsection}
\providecommand{\nobreakdash}{\nobreak\hbox{-}}
\setlength{\emergencystretch}{2em}
\multlinegap0pt
\usepackage{amssymb}
\usepackage{xfrac}
\usepackage{breqn}
\newtheorem{thm}{Theorem}
\title{Linear combination of bilateral gamma random variables: distributional theory and approximations}
\author{Kalyan Barman, Palaniappan Vellaisamy}
\date{Source: arXiv:2603.22112v1 (CC BY 4.0)}
\begin{document}
\maketitle

\noindent\emph{Source and licence.} Linear combination of bilateral gamma random variables: distributional theory and approximations. Version arXiv:2603.22112v1 (CC BY 4.0), distributed by arXiv under the Creative Commons Attribution 4.0 International licence, http://creativecommons.org/licenses/by/4.0/, as recorded in the arXiv licence metadata for this exact version and shown on the abstract page https://arxiv.org/abs/2603.22112v1. Full licence notice: \texttt{licenses/arxiv-2603.22112-CC-BY-4.0.txt}.\par\smallskip

\noindent\textit{Editorial scope.} Counted unit: the article's thm th3 (source0.tex lines 710--719). The article's complete original proof of it is delivered with it (source0.tex lines 720--772, one \texttt{proof} environment of 53 lines). No mathematics is written, repaired or re-derived: the statement and the proof are the article's own lines, in the article's own notation, and no retained line is substituted or re-flowed.  Retained uncounted as the source-local context the proof needs: lines 541--543; lines 630--632; lines 681--684.  Nothing was omitted from any retained span, so the package carries no omission record.  No retained line contains a citation, so the wrapper carries no bibliography and no bibliography entry.  No retained line contains a figure environment, an image-inclusion command or drawing code, so no image is delivered and the package declares no asset dependency.  Editorial formatting only, in addition: an AMS \texttt{amsart} preamble that restates the article's own declarations and macros used by the retained lines, namely the environments \texttt{thm} on separate counters, together with the standard packages those lines need.  The retained environments print in this document as Theorem 1 in order of appearance, whereas the article prints the same items with the numbering of its own full document.  The complete native source of the article as distributed in the arXiv e-print for this exact version is bundled unchanged as \texttt{sources/196998/source0.tex}.
\begin{align}\label{cumu}
	C_k(T_n)=(k-1)! \int_{\mathbb{{R}}}u^{k}\nu_{T_n}(du),~k\in\mathbb{N}.
\end{align} 

\begin{equation}\label{PP2:e15}
\mathcal{A}f(x)=-xf(x)+\displaystyle\int_{\mathbb{R}}f(x+u)u\nu_{T_n}(du)= h(x)-\mathbb{E}(h(T_n)),
\end{equation}

	\begin{equation}\label{PP2:SolSe}
	f_{h}(x):=-\displaystyle
	\int_{0}^{\infty}\frac{d}{dx}P_{t}h(x)dt,
	\end{equation}

\begin{thm}\label{th3}
	For $h\in\mathcal{W}_{k+1}$, let $f_{h}$ be defined in \eqref{PP2:SolSe}. Then
	\begin{align}\label{PP2:pr1}
	\|f_{h}^{(k)}\|\leq \frac{1}{k+1} ,~k\geq 0,
	\end{align}	
	where $f^{(k)}$, $k\geq 1$, denotes the $k$th derivative of $f$ with $f^{(0)}=f$. Also,
	\begin{align}
		|xf_{h}^{\prime}(x)| \leq 2+\frac{1}{2}|\mathbb{E}(T_n)| \text{ and }|xf_{h}^{\prime\prime}(x)| \leq 2+\frac{1}{3}|\mathbb{E}(T_n)|.
	\end{align}
\end{thm}

\begin{proof}
	For $h\in \mathcal{W}_{k+1}$,
	\begin{align*}
	\|f_h\|&=\sup_{x\in\mathbb{R}} \left|-\displaystyle
	\int_{0}^{\infty}\frac{d}{dx}P_{t}h(x)dt\right|\\
	&=\sup_{x\in\mathbb{R}} \left| -\displaystyle
	\int_{0}^{\infty}e^{-t}\int_{\mathbb{R}}h^{(1)}(xe^{-t}+y)F_{X_{(t)}}(dy)dt \right|\\
	&\leq \|h^{(1)}\| \left|\int_{0}^{\infty}e^{-t}dt \right|
	= \|h^{(1)}\|\leq 1.
	\end{align*}
	\noindent
	Since $f_{h}$ is $k$-times differentiable, we have 
	\begin{align*}
	\|f_h^{(1)}\|&=\sup_{x\in\mathbb{R}} \left| -\displaystyle
	\int_{0}^{\infty}e^{-2t}\int_{\mathbb{R}}h^{(2)}(xe^{-t}+y)F_{X_{(t)}}(dy)dt \right|\\
	&\leq \|h^{(2)}\| \left|\int_{0}^{\infty}e^{-2t}dt \right|\\	
	&= \frac{\|h^{(2)}\|}{2} \leq \frac{1}{2}.
	\end{align*}
	\noindent
	 Indeed, it follows by induction that $\|f_{h}^{(k)}\|\leq \frac{1}{k+1},~k\geq 1.$
	 
	
	
	
	\vskip 1ex
	\noindent
	Now differentiating both sides of \eqref{PP2:e15} gives
	\begin{align}\label{sed1}
		-xf^{\prime}(x)-f(x)+\int_{\mathbb{{R}}}f^{\prime}(x+u)u\nu_{T_n}(du)=h^{\prime}(x).
	\end{align}
	\noindent
	So,
	\begin{align*}
		|xf^{\prime}(x)|& \leq \|f\| +\|h^{\prime}\|+\left|\int_{\mathbb{{R}}}f^{\prime}(x+u)u\nu_{T_n}(du)\right|\\
		&\leq \left( 2 + \|f^\prime\|\left|\int_{\mathbb{{R}}}u\nu_{T_n}(du)\right| \right)\\
		& =2+\frac{1}{2}|C_1(T_n)| ~(\text{see } \eqref{cumu})\\
		&=2+\frac{1}{2}|\mathbb{E}(T_n)|.
	\end{align*}
	\noindent
	Again differentiating both sides of \eqref{sed1} gives
	\begin{align}
			-xf^{\prime\prime}(x)-2f^{\prime}(x)+\int_{\mathbb{{R}}}f^{\prime\prime}(x+u)u\nu_{T_n}(du)=h^{\prime\prime}(x).
	\end{align}
	\noindent
	So,
	\begin{align*}
		|xf^{\prime\prime}(x)| &\leq 2\|f^{\prime}\| +\|h^{\prime\prime}\|+ \|f^{\prime\prime}\|\left|\int_{\mathbb{{R}}}u\nu_{T_n}(du)\right|\\
		& \leq 2 +\frac{1}{3}|C_1(T_n)|\\
		&=2+\frac{1}{3}|\mathbb{E}(T_n)|.
	\end{align*}
	\noindent
	This proves the result.
\end{proof}

\end{document}
